We study nested Hermite interpolation with double nodes at the points $y_j=2\pi n_j$, where $n_j$ runs through the
prime powers, for three kernels: the Fourier transform of $\Xi^2$, which defines an explicit family of test functions
$\Xi(t)^2P_J(t^2)$ for the explicit formula of the Riemann zeta function, all vanishing at every non-trivial zero; its archimedean part; and the
kernel $e^{-y}$, which gives a polynomial model, equivalently a radial Fourier interpolation problem on the circles
$\lvert v\rvert=\sqrt{2n}$ in $\RR^2$. A kernel-free calculus (Christoffel, ladder, insertion and half-step identities)
reduces positivity of the interpolant beyond the nodes, in the model, to one likelihood-ratio law per level. A
computer-assisted certificate verifies this law at every level up to $160$, which proves the positivity of the front of
the model for every $J\le80$. We prove far-node asymptotics for the model and for the true kernel, show that sign
regularity of the associated five-function systems cannot give the tail of the model, and show that a natural sign
property fails at the critical node density. For the true family we reduce the open hypotheses of a criterion for
uniqueness and sharpness, proved in a companion paper, to explicit step laws, which remain open. Numerical observations
are collected in one section and are never used in proofs.
